\chapter{系统实现}

\section{实现语言与运行环境}

本实验实现使用 Windows PowerShell 脚本与模块组织代码。选择 PowerShell 的主要原因是 Windows 系统默认提供运行环境，并且能够调用 .NET 标准库中的大整数、随机数、哈希和对称加密组件。同一 \verb|Project| 工程在此基础上增加 Python 后端，用于调用真实 BLS12-381 pairing 库。系统源码主要位于：

\begin{itemize}
    \item \verb|Project/src/LsssAbe/LsssAbe.psm1|：核心模块。
    \item \verb|Project/bin/lsss-abe.ps1|：命令行入口。
    \item \verb|Project/tests/smoke.ps1|：烟雾测试。
    \item \verb|Project/src/LsssAbe/pairing_backend.py|：真实 BLS12-381 pairing 后端。
    \item \verb|Project/tests/smoke_pairing.ps1|：真实 pairing 版烟雾测试。
\end{itemize}

核心模块没有依赖第三方包，主要使用以下 .NET 类型：

\begin{table}[H]
\centering
\caption{实现中使用的主要 .NET 组件}
\begin{tabular}{p{0.35\textwidth} p{0.52\textwidth}}
\toprule
组件 & 用途 \\
\midrule
\verb|System.Numerics.BigInteger| & 表示模 $p$ 下的大整数和指数。\\
\verb|RandomNumberGenerator| & 生成密码学安全随机字节。\\
\verb|SHA256| & 将属性名和消息指数映射到固定长度摘要。\\
\verb|Aes| & 使用 AES-CBC 加密明文载荷。\\
\verb|HMACSHA256| & 对 IV 和密文进行完整性认证。\\
\bottomrule
\end{tabular}
\end{table}

真实 pairing 后端运行在 \verb|E:\WHU\Kinding_Plan\LSSS-ABE\env| 环境中，主要依赖如表 \ref{tab:python-pairing-deps} 所示。

\begin{table}[H]
\centering
\caption{真实 pairing 后端使用的 Python 组件}
\label{tab:python-pairing-deps}
\begin{tabular}{p{0.30\textwidth} p{0.58\textwidth}}
\toprule
组件 & 用途 \\
\midrule
\verb|py-ecc| & 提供 BLS12-381 的 $G_1/G_2/G_T/\mathbb{Z}_p$、标量乘法、群加法、哈希到 $G_1$ 和 \verb|pairing|。\\
\verb|pycryptodome| & 在真实 pairing 后端中实现 AES-CBC 加密、PKCS7 填充和 HMAC 所需基础能力。\\
\verb|secrets|、\verb|hashlib|、\verb|hmac| & 提供随机标量采样、SHA-256 摘要和认证标签比较。\\
\bottomrule
\end{tabular}
\end{table}

\section{模块结构}

核心模块可以分为六类函数：

\begin{enumerate}
    \item 数论工具：模运算、模逆、快速幂、向量加法、向量数乘和内积。
    \item 群运算适配器：封装随机标量、属性哈希、群加减乘和配对接口。
    \item 密码原语：随机数、哈希、AES-CBC-HMAC。
    \item 策略与 LSSS：策略树读取、随机维度统计、门索引分配、矩阵构造。
    \item 线性重构：从用户属性筛选行并用高斯消元求 $\omega$。
    \item CP-ABE 主算法：Setup、KeyGen、Encrypt、Decrypt。
\end{enumerate}

这样的分组虽然都位于同一个 PowerShell 模块文件中，但逻辑边界比较清晰。CLI 脚本只调用导出的主函数，不直接操作内部数学细节。

在实现时，核心模块没有进一步拆成多个文件，主要是为了降低实验复现时的路径复杂度。虽然文件较长，但函数命名保持了统一前缀 \verb|LsssAbe|，并按照“工具函数在前、算法函数在后、导出函数最后”的顺序组织。Setup、KeyGen、Encrypt 和 Decrypt 四个导出主函数集中呈现算法主流程，访问树到 LSSS 矩阵的转换函数则集中呈现矩阵构造逻辑。

这种组织方式也对应实验三课件~[1]中的算法分块：先准备底层群运算和随机数，再实现 KeyGen、Encrypt、Decrypt 等主算法。基础实现把真实群运算替换为模 $p$ 指数运算，把 Charm-Crypto 依赖替换为 PowerShell 与 .NET 标准库，从而使实验环境更容易复现；真实 pairing 路线则通过 Python 后端补回真实 pairing 运算，用于检查同一 LSSS 主线在真实曲线群元素上的运行情况。

\section{群运算适配器设计}

为使当前实验实现具备向真实双线性群迁移的边界，核心模块新增了群运算适配器层。现有默认适配器为 \verb|ExponentSimulation|，它仍然使用整数指数表示群元素：群乘法对应指数加法，配对对应指数乘法。因此，当前密文格式和实验行为保持不变，但 Setup、KeyGen、Encrypt 和 Decrypt 不再直接散落书写底层指数运算，而是通过适配器函数完成。

适配器层的核心接口包括：

\begin{itemize}
    \item \verb|New-LsssAbeGroupRandomScalar|：生成 $\mathbb{Z}_p$ 中的随机标量。
    \item \verb|New-LsssAbeGroupHashAttribute|：把属性字符串映射为当前指数抽象模型中的标量。
    \item \verb|Invoke-LsssAbeGroupAdd|、\verb|Invoke-LsssAbeGroupSub|、\verb|Invoke-LsssAbeGroupMul|：封装群指数模型下的加减乘。
    \item \verb|Invoke-LsssAbeGroupPair|：封装配对接口；在当前指数模型中返回两个指数的乘积。
\end{itemize}

关键代码节选如下：

\begin{lstlisting}[language=bash,caption={指数抽象群运算适配器关键代码}]
function New-LsssAbeExponentGroupAdapter {
    param([Parameter(Mandatory)][System.Numerics.BigInteger]$p)

    [pscustomobject]@{
        name = "ExponentSimulation"
        scalarModulus = $p
        supportsRealPairing = $false
    }
}

function Invoke-LsssAbeGroupPair {
    param(
        [Parameter(Mandatory)][object]$Group,
        [Parameter(Mandatory)][System.Numerics.BigInteger]$Left,
        [Parameter(Mandatory)][System.Numerics.BigInteger]$Right
    )

    Invoke-LsssAbeGroupMul -Group $Group -a $Left -b $Right
}
\end{lstlisting}

在当前 \verb|Project| 中，命令行入口新增 \verb|-Adapter pairing| 参数。当该参数被启用时，PowerShell 不再调用指数抽象模块中的四个主函数，而是把 setup、keygen、encrypt、decrypt 命令转发给 Python 后端 \verb|pairing_backend.py|。这种桥接方式保持了 CLI 使用习惯，同时避免在 PowerShell 中手动实现复杂的椭圆曲线群元素运算。

\begin{lstlisting}[language=bash,caption={CLI 中转发真实 pairing 后端的关键代码}]
if ($Adapter -eq "pairing") {
    switch ($Command) {
        "setup" {
            Invoke-PairingBackend -BackendArgs @("setup", "--out-dir", $OutDir)
            break
        }
        "keygen" {
            Invoke-PairingBackend -BackendArgs @("keygen", "--pub", $Pub,
                "--msk", $Msk, "--attrs", $Attrs, "--out-dir", $OutDir)
            break
        }
        "encrypt" {
            Invoke-PairingBackend -BackendArgs @("encrypt", "--pub", $Pub,
                "--policy", $Policy, "--input", $In, "--output", $Out)
            break
        }
        "decrypt" {
            Invoke-PairingBackend -BackendArgs @("decrypt", "--pub", $Pub,
                "--sk", $Sk, "--input", $In, "--output", $Out)
            break
        }
    }
    return
}
\end{lstlisting}

\subsection{真实 pairing 后端}

真实 pairing 后端的核心职责是把 CP-ABE 中的群元素计算放到 BLS12-381 上执行。后端从 \verb|py_ecc.optimized_bls12_381| 导入 $G_1$、$G_2$、$G_T$、曲线阶、标量乘法、群加法、逆元和配对函数，并从 \verb|py_ecc.bls.hash_to_curve| 导入 \verb|hash_to_G1|。关键导入如下：

\begin{lstlisting}[language=Python,caption={真实 pairing 后端关键导入}]
from py_ecc.bls.hash_to_curve import hash_to_G1
from py_ecc.optimized_bls12_381 import (
    G1, G2, Z1, Z2, FQ, FQ2, FQ12,
    add, multiply, neg, pairing,
    curve_order, is_inf, is_on_curve,
)
\end{lstlisting}

Setup 阶段在真实后端中选择随机标量 $\alpha$ 和 $a$，生成 $g_1^a$ 和 $e(g_2,g_1)^\alpha$。Encrypt 阶段为每个矩阵行生成 $C_i\in G_1$ 和 $D_i\in G_2$，Decrypt 阶段调用 \verb|pairing| 计算 $e(L,C_i)$、$e(D_i,K_{\rho(i)})$ 和 $e(C_0,K)$。核心代码如下：

\begin{lstlisting}[language=Python,caption={真实 pairing 后端的加解密核心片段}]
egg_alpha = pairing(G2, G1) ** alpha

ci = add(multiply(g1_a, lambdas[i]),
         neg(multiply(h_attr, blind[i])))
di = multiply(g2, blind[i])

term = pairing(L, Ci[row_index]) * pairing(Di[row_index], Kx[attr])
A *= term ** weights["omega"][pos]
masked = pairing(C0, K) * A.inv()
message_gt = C * masked.inv()
\end{lstlisting}

后端还实现了群元素 JSON 编码与反序列化检查。对于 $G_1$ 和 $G_2$ 点，程序会检查曲线成员关系，并用曲线阶乘法检查子群成员关系；对于 $G_T$ 元素，程序检查系数数量是否为 12。该处理使密钥和密文不再只是整数文件，而是包含真实曲线点和目标群元素的结构化 JSON。

\section{数据结构设计}

\subsection{访问控制树}

策略采用 JSON 表示。叶子节点包含 \verb|attr| 字段：

\begin{lstlisting}[caption={叶子属性节点示例}]
{ "attr": "A" }
\end{lstlisting}

内部节点包含阈值 \verb|k| 和子节点数组 \verb|children|：

\begin{lstlisting}[caption={阈值门节点示例}]
{
  "k": 2,
  "children": [
    { "attr": "A" },
    { "attr": "B" }
  ]
}
\end{lstlisting}

示例策略 $(A\wedge B)\vee C$ 写成：

\begin{lstlisting}[caption={示例访问策略 JSON}]
{
  "k": 1,
  "children": [
    {
      "k": 2,
      "children": [
        { "attr": "A" },
        { "attr": "B" }
      ]
    },
    { "attr": "C" }
  ]
}
\end{lstlisting}

策略对象进入矩阵构造前会先经过递归校验。该校验把属性叶子和阈值门分开处理，能够在加密前拒绝空属性、缺失阈值、空子节点、非法阈值以及叶子节点错误携带子节点等情况。关键代码节选如下：

\begin{lstlisting}[language=bash,caption={访问策略递归校验关键代码}]
function Assert-LsssAbePolicyNode {
    param([Parameter(Mandatory)][AllowNull()][object]$Node)

    if ($null -eq $Node) { throw "Invalid policy node: node must not be null." }

    $attr = Get-LsssAbeProp -Obj $Node -Name "attr"
    $childrenValue = Get-LsssAbeProp -Obj $Node -Name "children"

    if ($null -ne $attr) {
        if (-not ($attr -is [string])) { throw "Invalid policy leaf: attr must be a string." }
        if ([string]::IsNullOrWhiteSpace([string]$attr)) { throw "Invalid policy leaf: attr must not be empty." }
        if ($null -ne $childrenValue) { throw "Invalid policy node: leaf must not also define children." }
        return
    }

    $kValue = Get-LsssAbeProp -Obj $Node -Name "k"
    if ($null -eq $kValue) { throw "Invalid policy gate: missing k." }
    if ($null -eq $childrenValue) { throw "Invalid policy gate: missing children." }

    $children = @($childrenValue)
    if ($children.Length -lt 1) { throw "Invalid policy gate: children must not be empty." }
    if (-not ($kValue -is [byte] -or $kValue -is [int16] -or $kValue -is [int32] -or $kValue -is [int64])) {
        throw "Invalid policy gate: k must be an integer."
    }
    $k = [int]$kValue
    if ($k -lt 1) { throw "Invalid policy gate: k must be at least 1." }
    if ($k -gt $children.Length) { throw "Invalid policy gate: k must not exceed child count." }

    foreach ($c in $children) { Assert-LsssAbePolicyNode -Node $c }
}
\end{lstlisting}

这一段校验代码的意义在于把错误尽量前移。如果非法策略直接进入矩阵构造过程，后续错误可能表现为数组越界、空对象访问或线性代数维度不一致，定位成本较高。现在系统在加密前即可给出 \verb|Invalid policy| 相关错误，实验中的非法策略输入场景也正是验证这一点。

\subsection{LSSS 表示}

构造结果包含矩阵 $M$ 和映射 $\rho$。在内存中，矩阵行是 \verb|BigInteger[]|；写入 JSON 时，每个整数转为字符串，避免大整数在 JSON 解析时发生精度损失。

图 \ref{fig:policy-lsss-matrix} 展示了示例策略到矩阵的转换结果。

\begin{figure}[H]
\centering
\includegraphics[width=0.95\textwidth]{policy_lsss_matrix.png}
\caption{示例策略到 LSSS 矩阵的转换}
\label{fig:policy-lsss-matrix}
\end{figure}

图 \ref{fig:ct-json-lsss-screenshot} 展示了实际密文文件 \verb|ct.json| 中的 \verb|lsss| 字段。密文不仅保存对称加密载荷，也保存由访问控制树生成的矩阵 $M$ 和行属性映射 $\rho$，使矩阵结构能够与前文的策略转换结果相互对应。

\begin{figure}[H]
\centering
\includegraphics[width=0.85\textwidth]{ct_json_M_rho.png}
\caption{密文 ct.json 中 LSSS 矩阵与 $\rho$ 映射截图}
\label{fig:ct-json-lsss-screenshot}
\end{figure}

\subsection{密钥与密文}

公共参数、主密钥、用户私钥和密文都使用 JSON 存储。这样做的优点是便于人工检查和报告分析。例如密文中不仅包含 CP-ABE 分量，还包含策略树、LSSS 矩阵、属性映射和对称加密载荷。

JSON 存储还有一个现实好处：它让每个阶段的输出都能成为下一阶段的输入文件。Setup 输出的 \verb|public.json| 和 \verb|master.json| 被 KeyGen 和 Encrypt 使用；KeyGen 输出的 \verb|sk_*.json| 被 Decrypt 使用；Encrypt 输出的密文 JSON 同时携带策略和密文分量。这样整个实验过程可以被拆成一组清晰的命令，而不是隐藏在一个不可观察的内存流程中。

\begin{table}[H]
\centering
\caption{主要 JSON 对象字段}
\begin{tabular}{p{0.20\textwidth} p{0.25\textwidth} p{0.45\textwidth}}
\toprule
对象 & 字段 & 含义 \\
\midrule
公共参数 & \verb|p|、\verb|g|、\verb|ga|、\verb|egg_alpha| & 素数模数、生成元指数、公钥指数分量。\\
主密钥 & \verb|alpha|、\verb|a| & 主密钥指数。\\
用户私钥 & \verb|attrs|、\verb|K|、\verb|L|、\verb|Kx| & 用户属性集合和属性密钥分量。\\
密文 & \verb|policy|、\verb|lsss|、\verb|C|、\verb|C0|、\verb|Ci|、\verb|Di|、\verb|payload| & 策略、矩阵、CP-ABE 分量和对称密文载荷。\\
\bottomrule
\end{tabular}
\end{table}

真实 pairing 版沿用相同的高层字段名称，但字段内容不同。例如公共参数中的 \verb|g1_a| 保存 $G_1$ 点，\verb|egg_alpha| 保存 $G_T$ 元素；密文中的 \verb|C| 是 $G_T$ 元素，\verb|C0| 和 \verb|Di| 是 $G_2$ 点，\verb|Ci| 是 $G_1$ 点。因此，指数版和 pairing 版的 JSON 文件不能混用。保留相同字段语义的好处是报告中的算法结构仍然一致，而字段内部编码反映了底层群实现的差异。

\section{访问控制树到 LSSS 的实现}

实现树到矩阵转换的主函数是 \verb|ConvertTo-LsssAbeLsssFromPolicy|。其步骤如下：

\begin{enumerate}
    \item 调用 \verb|Get-LsssAbeGateRandomCount| 遍历策略树，计算所有阈值门需要的随机维度数量。
    \item 令矩阵列数为 $1+\sum(k-1)$，其中第一列对应秘密 $s$。
    \item 调用 \verb|New-LsssAbeGateIndexMap| 为每个阈值门分配随机变量列索引。
    \item 从根节点向下递归访问，根节点份额向量设为 $(1,0,\dots,0)$。
    \item 对每个内部阈值门，用父节点份额向量作为多项式常数项，用单位向量作为高阶系数。
    \item 对每个子节点编号 $1,\dots,d$ 求多项式值，得到子节点份额向量。
    \item 到达叶子时，将份额向量作为矩阵一行，并记录对应属性。
\end{enumerate}

该方法的核心优点是统一支持 AND、OR 和一般阈值门。AND 和 OR 不需要额外写分支逻辑，只需按照 $k$-of-$n$ 的通用规则处理。

实现中最容易出错的地方是份额向量维度和随机列索引的对应关系。为避免递归过程中动态扩展矩阵列数，程序先遍历整棵树计算总随机维度，再为每个阈值门提前分配列索引。这样每个叶子生成的行向量长度都相同，后续序列化和高斯消元也能使用统一维度。

Liu、Cao 和 Wong~[4]的阈值树转换算法使用“逐步插入”方式把非叶节点替换成相应阈值门的 LSSS 行。本项目没有采用同样的数据结构，但采用了相同的工程目标：输入是阈值访问树，输出是矩阵 $M$ 和行属性映射 $\rho$。递归传播份额向量的实现更适合 PowerShell 中直接展示每个叶子的最终行向量，因此也更便于在报告中解释矩阵来源。

对应实现中，\verb|ConvertTo-LsssAbeLsssFromPolicy| 先统计随机维度、分配列索引，再从根节点递归向下传播份额向量。内部节点使用父节点份额作为多项式常数项，使用新分配的单位向量作为高阶系数；叶子节点则直接写入矩阵行。核心递归代码节选如下：

\begin{lstlisting}[language=bash,caption={树到 LSSS 矩阵的递归构造代码}]
$rootShare = New-LsssAbeVecZeros -Len $n
$rootShare[0] = [System.Numerics.BigInteger]::One

function Visit {
    param([object]$Node, [System.Numerics.BigInteger[]]$ShareVec)

    $attr = Get-LsssAbeProp -Obj $Node -Name "attr"
    if ($null -ne $attr) {
        $rows.Add([pscustomobject]@{ attr = [string]$attr; row = $ShareVec })
        return
    }

    $k = [int](Get-LsssAbeProp -Obj $Node -Name "k")
    $children = @((Get-LsssAbeProp -Obj $Node -Name "children"))
    $gateKey = [System.Runtime.CompilerServices.RuntimeHelpers]::GetHashCode($Node)
    $gateRandIndices = @($indexMap[$gateKey])

    $coeffs = New-Object System.Collections.Generic.List[object]
    $coeffs.Add($ShareVec)
    for ($j = 1; $j -le ($k - 1); $j++) {
        $coeffs.Add((New-LsssAbeVecUnit -Len $n -Index $gateRandIndices[$j - 1] -p $p))
    }

    for ($i = 0; $i -lt $children.Length; $i++) {
        $t = [System.Numerics.BigInteger]($i + 1)
        $childShare = New-LsssAbeVecZeros -Len $n
        for ($j = 0; $j -le ($k - 1); $j++) {
            $pow = Invoke-LsssAbePowZp -x $t -e ([System.Numerics.BigInteger]$j) -p $p
            $term = Invoke-LsssAbeVecScale -v $coeffs[$j] -k $pow -p $p
            $childShare = Invoke-LsssAbeVecAdd -a $childShare -b $term -p $p
        }
        Visit -Node $children[$i] -ShareVec $childShare
    }
}

Visit -Node $PolicyTree -ShareVec $rootShare
\end{lstlisting}

高斯消元部分采用模 $p$ 运算，而不是普通有理数运算。每次找到主元后都要计算模逆并归一化主元行，再消去其他行中的同列元素。若消元后存在系数全为零但常数项非零的行，说明方程组不相容，也就是用户属性无法重构秘密。这个返回 \verb|$null| 的设计使解密函数可以用一个清晰的分支处理失败路径。

\section{模运算与高斯消元}

由于所有运算都在 $\mathbb{Z}_p$ 中进行，普通整数的负数取余和除法不能直接使用。模块中定义了模归一化函数，保证结果落在 $[0,p-1]$ 范围内。模逆通过费马小定理实现：

\[
    x^{-1}\equiv x^{p-2}\pmod p.
\]

重构系数求解由 \verb|Solve-LsssAbeWeights| 完成。它将问题写成：

\[
    M_I^T\omega=e_1.
\]

其中 $I$ 是用户属性能匹配的行集合。函数构造增广矩阵，并执行模 $p$ 下的高斯消元。如果消元后出现系数全为 $0$ 但常数项非 $0$ 的行，则方程无解，返回空结果。解密函数据此抛出不满足策略的错误。

行筛选和增广矩阵构造的关键代码如下。实现中没有对重复属性行进行去重：只要某一行的 $\rho(i)$ 出现在用户属性集合中，该行就会加入待求解系统，因此能够支持多行映射到同一属性的情况。

Waters~[3]特别指出，属性可能在访问矩阵中对应多行，也就是 $\rho$ 不一定是单射；这与访问树中同一个属性出现在多个叶子节点的情况相同。因此，代码在筛选行时不能把属性集合简单去重后只保留第一行，而必须保留所有满足 $\rho(i)\in S$ 的矩阵行。第五章的重复属性测试正是针对这一点设计的。

\begin{lstlisting}[language=bash,caption={重构系数求解中的行筛选与增广矩阵构造}]
$rowIndicesList = New-Object System.Collections.Generic.List[int]
for ($rowIdx = 0; $rowIdx -lt $rho.Length; $rowIdx++) {
    if ($Attrs -contains $rho[$rowIdx]) { $rowIndicesList.Add($rowIdx) }
}
if ($rowIndicesList.Count -eq 0) { return $null }

$n = $M[0].Length
$l = $rowIndicesList.Count
$A = New-Object 'System.Numerics.BigInteger[][]' $n
for ($r = 0; $r -lt $n; $r++) {
    $row = New-Object 'System.Numerics.BigInteger[]' ($l + 1)
    for ($c = 0; $c -lt $l; $c++) {
        $rowIndex = $rowIndicesList[$c]
        $row[$c] = Invoke-LsssAbeMod -x ($M[$rowIndex][$r]) -p $p
    }
    $row[$l] = if ($r -eq 0) { [System.Numerics.BigInteger]::One } else { [System.Numerics.BigInteger]::Zero }
    $A[$r] = $row
}
\end{lstlisting}

\section{加密载荷实现}

本项目的 CP-ABE 部分只负责保护随机消息指数 $m$，实际明文通过对称加密处理。对称加密流程如下：

\begin{enumerate}
    \item 将 $m$ 转为字符串并计算 SHA-256，得到基础密钥材料。
    \item 通过追加不同标签 \verb|AES| 和 \verb|HMAC|，派生加密密钥和认证密钥。
    \item 随机生成 16 字节 IV。
    \item 使用 AES-CBC 和 PKCS7 填充加密明文。
    \item 对 \verb|IV || ciphertext| 计算 HMAC-SHA256。
    \item 输出 Base64 编码的 \verb|nonce|、\verb|tag| 和 \verb|data|。
\end{enumerate}

解密时先重新计算 HMAC 并比较标签，只有标签正确才执行 AES 解密。这一设计使密文被篡改时能够尽早失败。

在 CP-ABE 指数抽象层，加密函数为每一行计算秘密份额 $\lambda_i$ 和盲化份额 $w_i$，再生成行密文分量 $C_i,D_i$。相关代码如下：

\begin{lstlisting}[language=bash,caption={加密阶段行分量生成关键代码}]
$group = Get-LsssAbeGroupAdapterFromPublic -Public $Public
$s = New-LsssAbeGroupRandomScalar -Group $group
$v[0] = $s
$u[0] = [System.Numerics.BigInteger]::Zero
for ($j = 1; $j -lt $n; $j++) {
    $v[$j] = New-LsssAbeGroupRandomScalar -Group $group
    $u[$j] = New-LsssAbeGroupRandomScalar -Group $group
}

for ($i = 0; $i -lt $l; $i++) {
    $lambda[$i] = Invoke-LsssAbeVecDot -a $M[$i] -b $v -p $p
    $w[$i] = Invoke-LsssAbeVecDot -a $M[$i] -b $u -p $p
}

for ($i = 0; $i -lt $l; $i++) {
    $attr = [string]$rho[$i]
    $h = New-LsssAbeGroupHashAttribute -Group $group -Text $attr
    $ciLeft = Invoke-LsssAbeGroupMul -Group $group -a $a -b $lambda[$i]
    $ciRight = Invoke-LsssAbeGroupMul -Group $group -a $h -b $w[$i]
    $ciExp = Invoke-LsssAbeGroupSub -Group $group -a $ciLeft -b $ciRight
    $diExp = Invoke-LsssAbeGroupMod -Group $group -x $w[$i]
    $Ci += $ciExp.ToString()
    $Di += $diExp.ToString()
}
\end{lstlisting}

这里把 CP-ABE 保护的对象设计为随机消息指数 $m$，而不是直接把文件字节放入群指数中。这样既符合混合加密的常见思路，也避免了大文件直接参与公钥运算的问题。实际明文由 AES-CBC-HMAC 处理，CP-ABE 只负责保护派生对称密钥所需的随机指数。

解密函数先调用 \verb|Solve-LsssAbeWeights| 求重构系数，再按 $\sum_i\omega_i(C_iL+D_iK_{\rho(i)})$ 计算抵消项。关键代码如下：

\begin{lstlisting}[language=bash,caption={解密阶段重构与恢复消息指数关键代码}]
$group = Get-LsssAbeGroupAdapterFromPublic -Public $Public
$weights = Solve-LsssAbeWeights -M $M -rho $rho -Attrs $attrs -p $p
if ($null -eq $weights) { throw "Attributes do not satisfy policy." }

$Aexp = [System.Numerics.BigInteger]::Zero
for ($pos = 0; $pos -lt $rowIndices.Length; $pos++) {
    $rowIndex = $rowIndices[$pos]
    $attr = [string]$rho[$rowIndex]
    $ciExp = [System.Numerics.BigInteger]::Parse($Ci[$rowIndex])
    $diExp = [System.Numerics.BigInteger]::Parse($Di[$rowIndex])
    $term1 = Invoke-LsssAbeGroupPair -Group $group -Left $ciExp -Right $keyL
    $term2 = Invoke-LsssAbeGroupPair -Group $group -Left $diExp -Right $Kx[$attr]
    $term = Invoke-LsssAbeGroupAdd -Group $group -a $term1 -b $term2
    $weighted = Invoke-LsssAbeGroupMul -Group $group -a $term -b $omega[$pos]
    $Aexp = Invoke-LsssAbeGroupAdd -Group $group -a $Aexp -b $weighted
}

$eC0K = Invoke-LsssAbeGroupPair -Group $group -Left $C0 -Right $keyK
$Bexp = Invoke-LsssAbeGroupSub -Group $group -a $eC0K -b $Aexp
$mExp = Invoke-LsssAbeGroupSub -Group $group -a $Cexp -b $Bexp
\end{lstlisting}

\section{CLI 实现}

命令行入口 \verb|lsss-abe.ps1| 负责读取参数并分派到核心模块函数。典型运行流程为：

\begin{lstlisting}[language=bash,caption={完整命令行流程}]
powershell -NoProfile -ExecutionPolicy Bypass -File .\bin\lsss-abe.ps1 setup -OutDir .\tests\artifacts
powershell -NoProfile -ExecutionPolicy Bypass -File .\bin\lsss-abe.ps1 keygen -Pub .\tests\artifacts\public.json -Msk .\tests\artifacts\master.json -OutDir .\tests\artifacts -Attrs "A,B"
powershell -NoProfile -ExecutionPolicy Bypass -File .\bin\lsss-abe.ps1 encrypt -Pub .\tests\artifacts\public.json -Policy .\examples\policy.json -In .\examples\message.txt -Out .\tests\artifacts\ct.json
powershell -NoProfile -ExecutionPolicy Bypass -File .\bin\lsss-abe.ps1 decrypt -Pub .\tests\artifacts\public.json -Sk .\tests\artifacts\sk_A_B.json -In .\tests\artifacts\ct.json -Out .\tests\artifacts\out.txt
\end{lstlisting}

CLI 的优点是阶段边界清晰，能够将每一步算法输出固定为可检查的文件。当前实现已经在属性名参与私钥文件名生成前执行安全化处理，例如 \verb|role:admin| 会生成 \verb|sk_role_admin.json|，同时私钥 JSON 内部仍保留原始属性名。

文件名安全化函数只作用于输出文件名，不修改私钥对象内部的属性语义。关键代码如下：

\begin{lstlisting}[language=bash,caption={属性名生成私钥文件名前的安全化处理}]
function ConvertTo-SafeFileNamePart {
    param([Parameter(Mandatory)][string]$Text)

    $safe = $Text -replace '[\\/:*?"<>|]', '_'
    $safe = $safe -replace '\s+', '_'
    $safe = $safe.Trim(" ._")
    if ($safe -eq "") { $safe = "attr" }
    $safe
}

$attrArr = $Attrs.Split(",") | ForEach-Object { $_.Trim() } | Where-Object { $_ -ne "" }
$skObj = New-LsssAbeKeyGen -Public $public -Master $master -Attrs $attrArr
$safeAttrs = $attrArr | ForEach-Object { ConvertTo-SafeFileNamePart -Text $_ }
$name = "sk_" + ($safeAttrs -join "_") + ".json"
\end{lstlisting}

这一处理解决了 Windows 文件名和属性字符串语义之间的冲突。例如属性 \verb|role:admin| 中的冒号不能直接出现在 Windows 文件名中，但私钥 JSON 内部仍必须保存原始属性名，否则解密时无法和策略中的 \verb|role:admin| 匹配。测试场景四同时检查了文件名安全化和解密正确性。

\section{实现边界}

本系统实现了 CP-ABE 的实验流程，并在同一 \verb|Project| 中完成真实 pairing 初步接入，但仍有以下边界：

\begin{itemize}
    \item 指数抽象模型不提供真实群元素不可区分性和困难问题安全性，主要用于解释和快速验证。
    \item 真实 pairing 版已经使用 BLS12-381 群元素和配对运算，但尚未给出与某一论文安全模型完全一致的形式化安全证明。
    \item 基础指数版模数约 127 比特，仅用于演示，不应视为真实安全参数。
    \item 策略输入当前为 JSON 访问树，尚未实现自然布尔表达式解析。
    \item 输入校验仍可继续增强，例如补充畸形 JSON、缺失字段和密文数组长度不一致等异常输入检查。
\end{itemize}

这些边界并不影响课程实验对算法流程的观察，但构成实验结论的适用范围。指数版用于清楚展示 LSSS-CP-ABE 的代数结构，pairing 版用于说明底层群运算已经可以替换为真实双线性群实现；二者共同支撑实验完成度，但仍不能被直接表述为经过形式化证明的可部署 ABE 系统。
